\chapter{Learning with Errors (LWE)}\label{ch:lwe}

\section{Auffrischung: Ring}
Ein \textbf{Ring} ist eine Menge $R$ mit Addition $+$ und Multiplikation $\cdot$. Man kann addieren, subtrahieren und multiplizieren (mit $0$, $1$ und Gegenelementen), und es gilt das Distributivgesetz $a \cdot (b + c) = a \cdot b + a \cdot c$. Division ist nicht immer möglich, sonst wäre der Ring ein Körper.

Beispiele sind die ganzen Zahlen $\mathbb{Z}$ und die ganzen Zahlen modulo $q$, $\mathbb{Z}_q = \mathbb{Z}/q\mathbb{Z}$. Letzteres funktioniert wie eine Uhr mit $q$ Stunden: In $\mathbb{Z}_7$ gilt zum Beispiel $5 + 4 = 2$. Man rechnet mit ganzen Zahlen und nimmt nach jeder Operation den Rest bei der Division durch $q$.

\section{Das LWE-Problem}
Ein lineares Gleichungssystem über $\mathbb{Z}_q$ hat die Form
\[
    \mathbf{A} \cdot \mathbf{s} = \mathbf{b},
\]
wobei $\mathbf{A} \in \mathbb{Z}_q^{n \times m}$, $\mathbf{s} \in \mathbb{Z}_q^{m}$ und $\mathbf{b} \in \mathbb{Z}_q^{n}$. Ein Beispiel:
\[
    \begin{pmatrix}
        10 & 3 & 5 & 1 \\
        4  & 1 & 1 & 2 \\
        \vdots & \vdots & \vdots & \vdots \\
        3  & 1 & 1 & 5
    \end{pmatrix}
    \cdot
    \begin{pmatrix} s_1 \\ s_2 \\ s_3 \\ s_4 \end{pmatrix}
    =
    \begin{pmatrix} 10 \\ 3 \\ \vdots \\ 8 \end{pmatrix}
\]
Ausgeschrieben sind das die Gleichungen
\begin{align*}
    10 \cdot s_1 + 3 \cdot s_2 + 5 \cdot s_3 + 1 \cdot s_4 &= 10 \\
    4 \cdot s_1 + 1 \cdot s_2 + 1 \cdot s_3 + 2 \cdot s_4 &= 3 \\
    &\;\;\vdots \\
    3 \cdot s_1 + 1 \cdot s_2 + 1 \cdot s_3 + 5 \cdot s_4 &= 8.
\end{align*}
Ein solches System ist mit dem Gauss-Algorithmus leicht lösbar. Fügt man jedoch kleine Fehlerwerte $\mathbf{e} \in \mathbb{Z}_q^{n}$ hinzu, erhält man
\[
    \mathbf{A} \cdot \mathbf{s} + \mathbf{e} = \mathbf{b},
\]
und das Lösen wird überraschend schwer (s. Verbindung zu Gitterproblemen).

\section{Decisional LWE (dLWE)}
Das LWE-Problem lässt sich auch als Entscheidungsproblem formulieren. Gegeben ist eine LWE-Aufgabe $(\mathbf{A}, \mathbf{b})$, wobei $\mathbf{s}$ und $\mathbf{e}$ geheim sind. Die Aufgabe ist zu erraten, ob $\mathbf{b}$ als $\mathbf{A} \cdot \mathbf{s} + \mathbf{e}$ berechnet oder zufällig gewählt wurde. Odded Regev hat gezeigt, dass beide Varianten gleich schwer sind.

Das ist für die Sicherheit eines Verschlüsselungssystems relevant: Es reicht nicht, dass ein Angreifer das Secret nicht findet.
Ein Verschlüsselte Text muss für einen Angreifer wie Zufall aussehen, egal ob eine $0$ oder eine $1$ verschlüsselt wurde.
Könnte er trotzdem das Bit erkennen, hätte er dLWE gelöst (Sicherheitsbeweis durch Reduktion).

\section{Verbindung zu Gitterproblemen}
\subsection{LWE als CVP}
Alle ganzzahligen Linearkombinationen der Spalten von $\mathbf{A}$ (modulo $q$) bilden ein Gitter. $\mathbf{A} \cdot \mathbf{s}$ ist ein Punkt dieses Gitters, und $\mathbf{b} = \mathbf{A} \cdot \mathbf{s} + \mathbf{e}$ liegt um den kleinen Fehler $\mathbf{e}$ daneben. Der Gitterpunkt, der $\mathbf{b}$ am nächsten liegt, ist deshalb fast immer $\mathbf{A} \cdot \mathbf{s}$, und daraus lässt sich $\mathbf{s}$ berechnen. Eine konkrete LWE-Instanz lässt sich also lösen, indem man das Closest Vector Problem löst.

\subsection{LWE als SVP}
Sei $\parallel$ die Konkatenation und $\mathbf{I}_n$ die $n \times n$-Einheitsmatrix. Wir betrachten das Gitter
\[
    \mathcal{L} = \{\mathbf{x} \in \mathbb{Z}^{m+n+1} \mid (\mathbf{A} \parallel \mathbf{I}_n \parallel (-\mathbf{b})) \cdot \mathbf{x} = \mathbf{0} \bmod q\}.
\]
Der Spaltenvektor $(\mathbf{s}, \mathbf{e}, 1)$ ist ein Element von $\mathcal{L}$, denn
\[
    \begin{pmatrix} \mathbf{A} & \mathbf{I}_n & -\mathbf{b} \end{pmatrix}
    \cdot
    \begin{pmatrix} \mathbf{s} \\ \mathbf{e} \\ 1 \end{pmatrix}
    = \mathbf{A} \cdot \mathbf{s} + \mathbf{e} - \mathbf{b} = \mathbf{b} - \mathbf{b} = \mathbf{0} \bmod q.
\]
$\mathbf{s}$ und $\mathbf{e}$ sind klein und die letzte Komponente ist $1$, der Vektor ist also sehr kurz. Die meisten anderen Gittervektoren haben Einträge in der Grössenordnung von $q$. $(\mathbf{s}, \mathbf{e}, 1)$ ist deshalb ein kürzester Vektor in $\mathcal{L}$, also eine Lösung des SVP. Wer ihn findet, erhält direkt das Secret $\mathbf{s}$ und den Fehlervektor $\mathbf{e}$ und löst damit das LWE-System.

